% Solutions to the exercises of Chapter 10.
\begin{solutions}
  \item With $F = \tfrac12r\T Wr$ and $W$ symmetric, a small change $\dd x$ changes the residual by $\dd r = J\,\dd x$ and the cost by
        $\dd F = r\T W\,\dd r = (J\T Wr)\T\dd x$, so $\nabla F = J\T Wr = \sum_k(Wr)_k\nabla r_k$, a combination of the residual gradients. Differentiating this sum
        by the product rule,
        \[
          \dd(\nabla F) = \sum_k\nabla r_k\,(W\,\dd r)_k + \sum_k(Wr)_k\,\nabla^2r_k\,\dd x = \Big(J\T WJ + \sum_k(Wr)_k\nabla^2r_k\Big)\dd x,
        \]
        which is \eqref{eq:nlls-derivatives}. The first term needs only first derivatives; the second is weighted by the residuals, which is why
        Gauss--Newton may drop it when the residuals are small.

  \item The implementation follows \cref{alg:gauss-newton} with $r(a, b) = ae^{-bt} - y$ and Jacobian columns $e^{-bt}$ and $-ate^{-bt}$; LM adds
        $\lambda\diag(J\T J)$ and adapts $\lambda$ (divide by 10 after a successful step, multiply by 10 after a failed one). On the data of
        \cref{fig:gn-lm} (25 samples on $[0, 1.5]$~s, $a = 2$, $b = 1.3$, noise $0.04$) both reach $(\hat a, \hat b) = (1.993, 1.293)$ from the figure's start.
        Starting instead from each point of a grid of $16\times21$ initial guesses with $a_0 \in [-1.75, 5.75]$ and $b_0 \in [-2, 8]$, LM converged from
        all 336 starts and Gauss--Newton from 237. Typical failures: from $(a_0, b_0) = (0.5, 6)$ the first Gauss--Newton step overshoots to a huge decay
        rate, after which $e^{-bt} \approx 0$ for all $t > 0$, $J$ becomes rank deficient and the iteration never recovers; from $(1, 4)$ it overflows. LM rejects
        such steps and shortens them until the cost decreases, which is what makes it the default choice when the initial guess is poor.

  \item Let $x^\star$ be a solution with $r(x^\star) = 0$ and $J^\star = J(x^\star)$ of full column rank, and let $e = x - x^\star$. Taylor expansion gives
        $r(x) = J^\star e + O(\|e\|^2)$ and $J(x) = J^\star + O(\|e\|)$, hence $r(x) = J(x)\,e + O(\|e\|^2)$. The Gauss--Newton step from $x$ is
        $\delta = -(J\T WJ)\inv J\T Wr(x) = -e + O(\|e\|^2)$, so the new error is $e + \delta = O(\|e\|^2)$: quadratic convergence. With a non-zero residual
        $r(x^\star)$ the omitted Hessian term $\sum_k(Wr)_k\nabla^2r_k$ is of the size of the residual, and the convergence is only linear, with a rate that
        grows with the residual. Fitting the exponential to noise-free data from $(2.3, 1.0)$ gives errors $8\times10^{-2}$, $3\times10^{-3}$, $3\times10^{-6}$,
        $4\times10^{-12}$: the number of correct digits doubles each time.

  \item With 100 points on $y = 0.5x + 1$ (noise $0.2$) of which 20 are replaced by uniform outliers, and IRLS using the normalized residuals
        $s_i = |r_i|/\hat\sigma$ with the robust scale $\hat\sigma = 1.4826\,\mathrm{median}|r_i - \mathrm{median}(r)|$:
        \begin{center}\small
        \begin{tabular}{@{}l c c@{}}
        \toprule
        method & slope & intercept \\
        \midrule
        least squares & 0.394 & 1.92 \\
        Huber IRLS ($c = 1.345$) & 0.489 & 1.04 \\
        Cauchy IRLS ($c = 2.385$) & 0.495 & 0.98 \\
        least squares on the inliers only & 0.496 & 0.98 \\
        \bottomrule
        \end{tabular}
        \end{center}
        Least squares is dragged far off by the outliers. Huber bounds their influence (each outlier still pulls with constant force $c$) and Cauchy
        nearly ignores them, matching the fit that knows which points are inliers. The robust scale estimate matters: with the scale computed by a
        plain standard deviation, which the outliers inflate, both kernels stop at slopes near $0.44$ and intercepts near $1.45$, because the
        outliers no longer look large enough to be down-weighted.

  \item Build a random problem with the bundle-adjustment structure: 8 ``cameras'' with 6 parameters and 60 ``landmarks'' with 3, each landmark
        observed by 3 random cameras, each observation contributing a $2\times6$ camera block and a $2\times3$ landmark block to $J$. Then the landmark block $C$
        of $M = J\T J$ is block diagonal, so $C\inv$ is computed block by block. Solving the reduced system
        $(A - BC\inv B\T)\,\delta c = g_c - BC\inv g_l$, then back-substituting $\delta l = C\inv(g_l - B\T\delta c)$, matches the direct solve of the full
        $228\times228$ system to $10^{-13}$. The reduced system is only $48\times48$; for real problems with thousands of landmarks the saving is enormous.

  \item Apply \cref{prop:block-inverse} to $M = \begin{bmatrix}A & B\\ B\T & C\end{bmatrix}$, eliminating the landmark block $C$ (the proposition with the roles of
        the blocks exchanged, as in Exercise 5.4): the camera block of $M\inv$ is $(M/C)\inv = (A - BC\inv B\T)\inv$. Hence the reduced camera matrix is the
        inverse of the camera block of $M\inv$ (numerically, to $5\times10^{-14}$ in the example above). Interpretation: at the optimum $M = J\T WJ$ is the
        information matrix of the estimate and $M\inv$ its covariance. The camera block of the covariance is the marginal covariance of the cameras,
        and its inverse is their marginal information. The reduced camera system is therefore exactly the information about the cameras after the
        landmarks have been marginalized out, the same operation that a sliding-window smoother performs when it marginalizes old states.

  \item Iteration $i$ linearizes $h$ at the current iterate $x_i$ (with $x_0 = \mu$), $H_i = \partial h/\partial x|_{x_i}$, and takes a Gauss--Newton step on the MAP
        cost from $x_i$. Its normal equations are
        \[
          \big(P\inv + H_i\T R\inv H_i\big)\,\delta = -P\inv(x_i - \mu) + H_i\T R\inv\big(z - h(x_i)\big).
        \]
        Adding $(P\inv + H_i\T R\inv H_i)(x_i - \mu)$ to both sides gives $(P\inv + H_i\T R\inv H_i)(x_{i+1} - \mu) = H_i\T R\inv\big(z - h(x_i) - H_i(\mu - x_i)\big)$, i.e.
        \[
          x_{i+1} = \mu + K_i\big(z - h(x_i) - H_i(\mu - x_i)\big), \qquad K_i = PH_i\T\big(H_iPH_i\T + R\big)\inv .
        \]
        For $i = 0$, $x_0 = \mu$ and the bracket reduces to $z - h(\mu)$: the first iteration is exactly the EKF update of \cref{prop:ekf-gn}. Later
        iterations re-linearize the measurement at the improved estimate; the term $-H_i(\mu - x_i)$ keeps the prior anchored at $\mu$, so the iteration
        minimizes the full MAP cost rather than repeatedly applying the measurement. After the last iteration the covariance is $(I - K_iH_i)P$.
\end{solutions}
